\chapter{측정 범위와 보고서 읽기}

\section{먼저 같은 것을 비교하고 있는지 확인한다}

``이 설계의 area''는 top-level을 말하지 않으면 불완전한 문장이다. 이 프로젝트는
다음 세 범위를 분리한다.

\begin{enumerate}
  \item 완전한 AXI4-Stream/AXI DMA/SmartConnect/Zynq PS 시스템
  \item 완전한 AXI BRAM Controller/SmartConnect/Zynq PS 시스템
  \item standalone \rtlfile{MDL} RTL의 out-of-context 구현
\end{enumerate}

첫 두 범위에는 processor interface와 data-movement IP가 들어간다. 세 번째는
가속기 RTL의 내부 경로와 자원을 보는 범위다. standalone 결과를 complete-system
결과와 나란히 놓을 수는 있지만, 서로 빼서 ``AXI 비용''이라고 단정해서는 안
된다. hierarchy boundary와 최적화 조건이 다르기 때문이다.

\section{재현 조건 표}

PPA 결과를 기록할 때 적어도 다음 항목을 표 위에 둔다.

\begin{longtable}{L{43mm}L{119mm}}
\toprule
항목 & 기록 이유 \\
\midrule
Tool/version & inference와 physical optimization이 버전에 따라 달라진다. \\
Device/package/speed grade & 같은 FPGA family라도 pinout과 배치 가능 영역, 속도가 달라진다. \\
Top-level/hierarchy & interface IP와 debug logic 포함 범위를 결정한다. \\
Clock constraint & utilization과 timing은 목표 주기에 반응한다. \\
Synthesis와 implementation directive & retiming, replication, placement 결과에 영향을 준다. \\
Post-synthesis와 post-route & 전자는 구조 확인, 후자는 실제 routing을 포함한 최종 판단에 쓴다. \\
\bottomrule
\end{longtable}

최종 세 범위는 Vivado 2020.2, \rtlfile{xc7z020clg484-1}, 6.666~ns
constraint로 맞췄다. 이 조건이 일치하므로 interface별 complete-system 결과를
비교할 수 있다.

\section{Timing report의 읽는 순서}

\begin{enumerate}
  \item \textbf{Clock/path group}: 의도한 clock의 synchronous path인지 본다.
  \item \textbf{Source/destination}: RTL의 어느 경계인지 찾는다.
  \item \textbf{WNS/TNS/failing endpoints}: 최악 정도와 문제의 폭을 구분한다.
  \item \textbf{Logic levels}: LUT와 CARRY4가 얼마나 이어지는지 본다.
  \item \textbf{Logic/route 비율}: Boolean 재작성과 register/placement 중 우선순위를 정한다.
  \item \textbf{Fanout}: 긴 route가 한 신호의 넓은 분배 때문인지 확인한다.
\end{enumerate}

\wns가 양수이면 해당 제약을 만족한다. TNS가 0이고 failing endpoint가 0인지도
함께 확인한다. 다음 식으로 slack에서 대략적인 주파수를 계산할 수 있다.
\[
T_{achieved}\approx T_{constraint}-\wns,\qquad
f_{approx}\approx \frac{1}{T_{achieved}}.
\]
이 값은 현재 placement에서의 근사치다. 다른 constraint로 다시 구현했을 때
같은 placement가 유지된다는 보장은 없다.

\section{실패한 synthesis도 데이터다}

초기 baseline synthesis에서는 coefficient RAM이 BRAM으로 inference되지 않고
register로 풀렸다. 결과는 LUT 563,561개, FF 191,202개, BRAM 0개, \wns
$-9.387$~ns였다. 이 숫자를 정상 architecture의 PPA로 사용하면 안 된다.
그러나 실패 원인을 찾는 데는 매우 강한 증거다. 보고서의
\texttt{RAM ``mem\_reg'' dissolved into registers}와 write-process 경고가 먼저
수정해야 할 대상이 RAM coding template임을 알려 주었다.

이후 inference가 정상화된 100~MHz routed baseline은 LUT 5,278개, FF 3,072개,
BRAM tile 7.5개, DSP48E1 4개가 되었다. area 최적화는 항상 이런 ``정상 mapping
baseline''을 확보한 뒤 시작해야 한다.

\begin{warningbox}[흔한 오판]
BRAM inference 실패로 수십만 LUT가 나온 run과 정상 run의 감소율을 설계 개선율로
쓰지 않는다. 전자는 기능 구조가 물리 자원에 잘못 mapping된 진단 결과다.
\end{warningbox}

\section{Hierarchy report를 이용한 원인 분리}

top utilization 한 줄만 보면 증가한 자원이 core 때문인지 interface 때문인지
알 수 없다. 최종 AXI4-Stream complete system 안의 \rtlfile{mdl} hierarchy는
4,739 LUT와 2,792 FF였고, standalone \rtlfile{MDL}은 4,736 LUT와 2,792 FF였다.
3-LUT 차이는 IP 경계를 넘는 최적화 수준의 작은 차이다. 이 일치는 standalone
결과가 통합 시스템 속 가속기와 같은 RTL을 대표한다는 확인 자료가 된다.

반면 complete-system BRAM은 stream 시스템 9 tile, BRAM-controller 시스템
6 tile이다. 이 차이는 arithmetic core가 아니라 DMA payload FIFO를 포함하는
interface 구성에서 나온다. hierarchy report는 이런 귀속을 가능하게 한다.
